\section{OS-World 的执行机制：观察、动作与交互循环}

\subsection{观察空间设计}
课程讲解了多源观察输入：屏幕截图、可访问性树（accessibility tree）、终端输出和自定义状态流。不同 Agent 可按需求选择视觉主导或文本主导的输入组合。

\begin{knowledgebox}{为什么截图 + 可访问性树要同时保留}
截图提供整体视觉上下文，但结构化程度低；可访问性树具备对象层级与属性信息，便于精准定位和可解释操作。二者互补能显著提高动作生成稳定性。
\end{knowledgebox}

\begin{figure}[H]
\centering
\includegraphics[width=0.88\textwidth]{frame-03-observation.jpg}
\caption{视频约 00:16:30：讲者展示 screenshot 与 accessibility tree 的组合观察形式}
\end{figure}

\subsection{动作空间与执行接口}
动作由可执行指令表达，覆盖点击、输入、滚动、快捷键与组合操作。Agent 并非直接输出 ``建议''，而是生成可在虚拟机中执行的行为代码。

\begin{importantbox}{动作格式标准化的价值}
统一动作格式能同时服务三件事：训练数据对齐、执行器稳定运行、评测脚本可复现。没有标准动作层，跨模型比较和大规模数据积累都会失败。
\end{importantbox}

\subsection{交互循环}
系统循环为 ``接收观察 $\rightarrow$ 规划动作 $\rightarrow$ 环境执行 $\rightarrow$ 新观察''。任务在达到目标、触发失败条件或到达最大步数时结束。这个循环把 Agent 训练目标从单步预测扩展为序列决策。

\begin{lstlisting}[caption={多模态 Agent 在 OS-World 的简化交互循环}]
obs = env.reset(task_config)
for step in range(max_steps):
    action = agent.predict(obs, instruction)
    obs, reward, done, info = env.step(action)
    if done:
        break
\end{lstlisting}

\subsection{工程细节：安全与可复现}
课程强调了虚拟机隔离和任务配置文件的重要性。研究环境必须可重置、可记录、可重放，这样才能做稳定对比实验并诊断失败轨迹。

\begin{warningbox}{常见失败点}
真实环境里最易被忽略的是 ``状态漂移''：缓存、登录态、弹窗与后台进程导致同任务在不同运行中行为不一致。必须通过环境快照、初始化脚本和日志追踪降低漂移影响。
\end{warningbox}

\subsection{本章小结}
OS-World 的核心不是界面截图本身，而是 ``结构化观察 + 标准动作 + 可重放执行'' 的闭环设计。这让多模态 Agent 研究进入可工程化迭代阶段。

